\chapter{Vacuum, Sectors, Energy, and the Budget Space}
\label{ch:vacuum_budget}

\section{Introduction and the Energy Budget Philosophy}

In continuum quantum field theory, the Dirac sea has infinite negative depth. Attempting to define particle-hole operators, charges, and Hamiltonians directly in the continuum leads immediately to divergent vacuum expectations and unbounded operators. 

To develop the finite-lattice bosonization without topological approximations or divergent series, we introduce the \emph{Energy Budget regularization scheme}. The entire state space is decomposed into finite-dimensional subspaces governed by two discrete parameters:
\begin{enumerate}
    \item Net charge sector deviation $N \in \mathbb{Z}$.
    \item Maximum excitation energy cutoff $K \in \mathbb{N}$.
\end{enumerate}

Throughout this chapter, let $L = 2h$ with $h \ge 1$ be an even positive integer. The momentum band is $\Lambda^* = \{-h+1, \dots, h\}$, and $\mathcal{F}(L)$ is the $2^L$-dimensional Fock space.

\section{Admissible Charge Sectors and Sector Grounds}

Because the single-particle band $\Lambda^*$ contains exactly $2h$ modes, an occupation configuration $S \subseteq \Lambda^*$ can have particle number $\#S$ between $0$ and $2h$. 

\begin{definition}[Admissible Charge Sectors]
The net charge deviation $N \in \mathbb{Z}$ of an occupation configuration $S \subseteq \Lambda^*$ relative to half-filling is defined by:
\begin{equation}
N(S) \coloneqq \#S - h \in \mathbb{Z}.
\end{equation}
A charge sector $N$ is \emph{admissible} if it can be realized by a subset of $\Lambda^*$:
\begin{equation}
-h \le N \le h.
\end{equation}
For any admissible sector $N$, the unique \emph{sector ground configuration} $\mathcal{S}_N \subset \Lambda^*$ is the set of the lowest $h + N$ momentum modes:
\begin{equation}
\mathcal{S}_N \coloneqq \{ k \in \Lambda^* \mid k \le N \} = \{-h+1, -h+2, \dots, N\}.
\end{equation}
The \emph{sector ground state} $|N\rangle_0 \in \mathcal{F}(L)$ is the corresponding Fock basis vector:
\begin{equation}
|N\rangle_0 \coloneqq |\mathcal{S}_N\rangle.
\end{equation}
When $N = 0$, $\mathcal{S}_0 = \{-h+1, \dots, 0\}$ is the reference half-filled Dirac vacuum, and $|0\rangle_0 = |\Omega\rangle$.
\end{definition}

\begin{lemma}[Cardinality of Sector Grounds]
For any admissible charge $N \in \{-h, \dots, h\}$:
\begin{equation}
\#\mathcal{S}_N = h + N, \qquad N(\mathcal{S}_N) = N.
\end{equation}
\end{lemma}

\begin{proof}
The modes in $\mathcal{S}_N$ are consecutive integers from $-h+1$ to $N$. The number of elements is $N - (-h+1) + 1 = N + h$. Subtracting $h$ yields $N(\mathcal{S}_N) = (h+N) - h = N$.
\end{proof}

\section{Integer Energy Grading and Observables}

To maintain strict mathematical rigor, all combinatorial quantities are defined as exact signed integers in $\mathbb{Z}$.

\begin{definition}[Integer Charge, Momentum, and Ground Energy Offset]
For any configuration $S \subseteq \Lambda^*$:
\begin{enumerate}
    \item The half-filled reference sea has integer momentum energy $E_\Omega$; the total bare momentum relative to it is:
    \begin{equation}
    E_\Omega \coloneqq \sum_{k\in\mathcal{S}_0}k=-\frac{h(h-1)}2,\qquad
    P(S) \coloneqq \sum_{k\in S}k-E_\Omega\in\mathbb Z.
    \end{equation}
    \item The \emph{sector ground reference energy offset} $E_0(N)$ is defined for $N \in \mathbb{Z}$ by:
    \begin{equation}
    E_0(N) \coloneqq \frac{N(N+1)}{2} \in \mathbb{Z}.
    \end{equation}
    Because one of $N$ or $N+1$ is always even, the product $N(N+1)$ is divisible by $2$ in $\mathbb{Z}$, so $E_0(N)$ is an exact integer for every $N \in \mathbb{Z}$.
    \item The \emph{excitation energy} $e(S)$ of configuration $S$ is defined by:
    \begin{equation}
    e(S) \coloneqq P(S) - E_0(N(S)).
    \end{equation}
\end{enumerate}
\end{definition}

\begin{quote}
\item[]\paragraph{Physical Intuition for the Ground Energy Offset $E_0(N)$.}
Why does the ground state of charge sector $N$ carry energy $E_0(N) = N(N+1)/2$?
Consider adding $N > 0$ particles to the half-filled Dirac sea $\mathcal{S}_0 = \{-h+1, \dots, 0\}$. To minimize total kinetic energy, the particles must occupy the lowest available empty modes immediately above the Fermi surface, which are the consecutive positive modes $k = 1, 2, \dots, N$. The total momentum added is:
\begin{equation}
\sum_{k=1}^N k = 1 + 2 + \dots + N = \frac{N(N+1)}{2}.
\end{equation}
Similarly, if $N < 0$ particles are removed, the lowest-energy configuration removes the highest occupied modes $k = 0, -1, \dots, N+1$ from the sea. The net momentum change is $-\sum_{k=N+1}^0 k = \frac{|N|(|N|-1)}{2} + |N| = \frac{N(N+1)}{2}$.
In 1D physics, this quadratic term $E_0(N) = \frac{1}{2} N^2 + \frac{1}{2} N$ reflects the macroscopic charging energy of a 1D Luttinger liquid: adding particles increases the Fermi momentum $k_F \propto N$, so the total kinetic energy increases quadratically as $\int_0^N k \, dk \sim \frac{1}{2} N^2$.
\end{quote}

\begin{definition}[Operator Observables]
The relative charge operator $N$ and excitation energy operator $E$ belong to $\operatorname{End}_{\mathbb C}(\mathcal F(L))$ and are defined on basis states by:
\begin{align}
N |S\rangle &\coloneqq N(S) |S\rangle = (\#S - h) |S\rangle, \\
E |S\rangle &\coloneqq e(S) |S\rangle.
\end{align}
\end{definition}

\begin{lemma}[Vacuum Ground State Evaluation]
The ground states satisfy:
\begin{equation}
P(\mathcal{S}_N) = E_0(N), \qquad e(\mathcal{S}_N) = 0 \quad \text{for all } N \in \{-h, \dots, h\}.
\end{equation}
In particular, $E |N\rangle_0 = 0$.
\end{lemma}

\begin{proof}
For $N\ge0$, summing the added modes gives
\begin{equation}
P(\mathcal S_N)=\sum_{k=1}^{N}k=\frac{N(N+1)}2=E_0(N).
\end{equation}
For $N<0$, reindexing the removed modes gives
\begin{equation}
P(\mathcal S_N)=-\sum_{j=0}^{-N-1}(N+1+j)=\frac{N(N+1)}2=E_0(N).
\end{equation}
Thus $e(\mathcal S_N)=P(\mathcal S_N)-E_0(N)=0$ in either case.
\end{proof}

\section{The Rank Formula and Ground State Uniqueness}

A central theorem of the energy budget is that the excitation energy $e(S)$ is strictly non-negative for all configurations, and vanishes if and only if $S$ is the sector ground configuration $\mathcal{S}_N$.

\begin{definition}[Ordered Modes, Ground Positions, and Particle Displacements]
Let $S \subseteq \Lambda^*$ be an occupation configuration with cardinality $\#S = h + N$.
\begin{enumerate}
    \item The \emph{ordered occupied modes} of $S$ are the elements of $S$ written in strictly increasing order:
    \begin{equation}
    s_0 < s_1 < s_2 < \dots < s_{h+N-1}.
    \end{equation}
    \item The \emph{ground reference positions} are the consecutive elements of the sector ground $\mathcal{S}_N$:
    \begin{equation}
    g_i \coloneqq -h + 1 + i \quad \text{for } i \in \{0, 1, \dots, h+N-1\}.
    \end{equation}
    \item For each particle index $i \in \{0, 1, \dots, h+N-1\}$, the \emph{particle displacement} $d_i$ is defined by:
    \begin{equation}
    d_i \coloneqq s_i - g_i \in \mathbb{Z}.
    \end{equation}
\end{enumerate}
\end{definition}

\begin{quote}
\item[]\paragraph{Physical Intuition for Particle Displacements.}
The ground state $\mathcal{S}_N$ is a compact Fermi droplet where particles occupy consecutive lowest modes $g_i$. Any excited state $S$ is created by displacing particles upward from their ground positions. The integer $d_i = s_i - g_i$ measures how far the $i$-th fermion has been promoted above its ground state reference position.
\end{quote}

\begin{theorem}[Rank Formula and Non-Negativity of Excitation Energy]
\label{thm:rank_formula}
Let $S \subseteq \Lambda^*$ be any configuration with cardinality $\#S = h + N$ (so $N(S) = N$). 
\begin{enumerate}
    \item The displacements are non-negative, non-decreasing, and bounded above by the band edge:
    \begin{equation}
    h+N=0\;:\;S=\emptyset,\qquad
    h+N>0\;:\;0\le d_0\le\cdots\le d_{h+N-1}\le 2h-(h+N)=h-N.
    \end{equation}
    \item The excitation energy is the sum of particle displacements:
    \begin{equation}
    e(S) = \sum_{i=0}^{h+N-1} d_i.
    \end{equation}
    \item Non-negativity: $e(S) \ge 0$ for all $S \subseteq \Lambda^*$.
    \item Uniqueness: $e(S) = 0$ if and only if $S = \mathcal{S}_N$.
\end{enumerate}
\end{theorem}

\begin{proof}
We prove each assertion:
\begin{enumerate}
    \item If $S=\emptyset$, the displacement assertions are vacuous. Otherwise $s_0\ge g_0$, so $d_0\ge0$; for consecutive indices, $s_{i+1}-s_i\ge1$ and $g_{i+1}-g_i=1$, hence $d_{i+1}-d_i\ge0$. Finally, $s_{h+N-1}\le h$ and $g_{h+N-1}=N$ give $d_{h+N-1}\le h-N$.

    \item Using $P(S)=\sum_{k\in S}k-E_\Omega$ and $P(\mathcal S_N)=E_0(N)$, the sorted occupations and ground positions give the telescoping identity
    \begin{equation}
    e(S)=P(S)-E_0(N)=\sum_{k\in S}k-\sum_{k\in\mathcal S_N}k
    =\sum_{i=0}^{h+N-1}(s_i-g_i)=\sum_{i=0}^{h+N-1}d_i.
    \end{equation}
    \item Since each $d_i \ge 0$, the sum of non-negative integers is non-negative: $e(S) \ge 0$.
    
    \item Since $d_i \ge 0$, the sum $\sum d_i = 0$ holds if and only if $d_i = 0$ for every $i \in \{0, \dots, h+N-1\}$.
    This means $s_i = g_i$ for all $i$, which implies $S = \mathcal{S}_N$.
\end{enumerate}
\end{proof}

\begin{corollary}[One-Dimensional Ground Subspace]
\label{cor:ground_subspace}
In each admissible sector $N$, the zero-excitation energy subspace is strictly one-dimensional:
\begin{equation}
\mathcal{H}(N, 0) \coloneqq \mathrm{span}_{\mathbb{C}} \{ |S\rangle \mid N(S) = N, \; e(S) = 0 \} = \mathbb{C} \cdot |N\rangle_0.
\end{equation}
\end{corollary}

\begin{quote}
\item[]\paragraph{Ground-state uniqueness and the gapless spectrum.}
In condensed matter physics, Corollary~\ref{cor:ground_subspace} carries profound physical meaning:
\begin{enumerate}
    \item \textbf{Non-Degenerate Ground State:} For any fixed particle number $h + N$, the ground state is unique (non-degenerate). There are no spontaneous symmetry-breaking degeneracies at zero temperature in 1D.
    \item \textbf{The Energy Gap:} The lowest possible non-zero excitation energy is obtained by setting $d_{h+N-1} = 1$ and all other $d_i = 0$, which corresponds to promoting the highest fermion across the Fermi level into mode $N+1$. The excitation energy of this state is $e(S) = 1$. In physical energy units, the excitation gap is $\Delta \epsilon = \frac{2\pi v_F}{L}$. As $L \to \infty$, the gap vanishes as $\mathcal{O}(1/L)$, confirming that the 1D free Fermi gas is a \emph{gapless (critical) system} with a continuous low-energy spectrum.
\end{enumerate}
\end{quote}

\section{The Frozen Margins Theorem}

Because the displacements $d_i$ are non-negative and non-decreasing, an upper bound on the total excitation energy $e(S) \le K$ strongly constrains the positions of particles at the bottom and top of the band.

\begin{theorem}[Frozen Margins]
\label{thm:frozen_margins}
Let $S \subseteq \Lambda^*$ be a configuration with $N(S) = N$ and excitation energy $e(S) \le K$.
\begin{enumerate}
    \item \textbf{Deep Sea Rigidity:} Every momentum mode deep below the Fermi surface is rigidly occupied:
    \begin{equation}
    \text{For all } k \in \Lambda^* \text{ with } k \le N - K, \quad k \in S.
    \end{equation}
    \item \textbf{High Momentum Rigidity:} Every momentum mode high above the Fermi surface is rigidly empty:
    \begin{equation}
    \text{For all } k \in \Lambda^* \text{ with } k > N + K, \quad k \notin S.
    \end{equation}
\end{enumerate}
\end{theorem}

\begin{proof}
By Theorem~\ref{thm:rank_formula}, $e(S) = \sum_{i=0}^{h+N-1} d_i \le K$. Since $d_i \ge 0$ for all $i$, each individual displacement satisfies $d_i \le K$.
\begin{enumerate}
    \item If a mode $k\le N-K$ were unoccupied, the resulting hole would force total displacement at least $N-k+1\ge K+1$, contradicting $\sum_i d_i=e(S)\le K$. Thus every such mode is occupied.
    \item If an occupied mode $k>N+K$ existed, the highest occupied mode would satisfy $s_{h+N-1}\ge k$ while its ground position is $g_{h+N-1}=N$. Hence $d_{h+N-1}\ge k-N>K$, contradicting $\sum_i d_i\le K$. Thus every such mode is empty.
\end{enumerate}
\end{proof}

\section{Energy Budget Subspaces and Filtered Operators}

\subsection{Coordinate Subspaces and Projector Algebra}
For $U\subseteq\mathcal P(\Lambda^*)$, define the coordinate subspace and its orthogonal projector by
\begin{equation}
\operatorname{coord}(U)=\operatorname{span}_{\mathbb C}\{|S\rangle:S\in U\},\qquad
P_U|S\rangle=\begin{cases}|S\rangle,&S\in U,\\0,&S\notin U,\end{cases}
\quad P_U\in\operatorname{End}_{\mathbb C}(\mathcal F(L)).
\end{equation}
Their algebra is
\begin{align}
P_U^2&=P_U, & P_U^\dagger&=P_U, & \operatorname{Ran}(P_U)&=\operatorname{coord}(U),\\
P_UP_V&=P_{U\cap V}, & P_{U\cup V}&=P_U+P_V-P_{U\cap V}.
\end{align}
Thus disjoint coordinate sets give orthogonal subspaces and $P_UP_V=0$. These identities follow by evaluating both sides on each occupation basis ket.


\begin{definition}[Energy Budget Subspaces]
\label{def:budget_subspace}
For sector $N \in \{-h, \dots, h\}$ and cutoff $K \in \mathbb{N}$:
\begin{enumerate}
    \item The \emph{fixed-energy subspace} $\mathcal{H}(N, E)$ is the span of basis states with charge $N$ and energy $E$:
    \begin{equation}
    \mathcal{H}(N, E) \coloneqq \mathrm{span}_{\mathbb{C}} \{ |S\rangle \mid N(S) = N, \; e(S) = E \}.
    \end{equation}
    \item The \emph{fixed-charge energy budget subspace} $B(N, K)$ is the internal orthogonal direct sum:
    \begin{equation}
    B(N, K) \coloneqq \bigoplus_{E=0}^K \mathcal{H}(N, E) = \mathrm{span}_{\mathbb{C}} \{ |S\rangle \mid N(S) = N, \; e(S) \le K \}.
    \end{equation}
    \item The \emph{charge-box budget subspace} $\mathcal{B}_{K, N_{\max}}$ for $N_{\max} \ge 0$ is:
    \begin{equation}
    \mathcal{B}_{K, N_{\max}} \coloneqq \bigoplus_{|N| \le N_{\max}} B(N, K) = \mathrm{span}_{\mathbb{C}} \{ |S\rangle \mid e(S) \le K, \; |N(S)| \le N_{\max} \}.
    \end{equation}
\end{enumerate}
The orthogonal coordinate projectors onto these subspaces are denoted $P_{N, K}$ and $P_{K, N_{\max}}$, respectively.
\end{definition}

\begin{definition}[Filtered Operator Parameters: Charge Shift and Energy Increase]
Let $A \in \mathrm{End}_{\mathbb{C}}(\mathcal{F}(L))$ be a linear operator.
\begin{enumerate}
    \item $A$ carries a \emph{charge shift} $q \in \mathbb{Z}$ if for every basis state $|S\rangle$, every output state $|S'\rangle$ in the support of $A |S\rangle$ has charge:
    \begin{equation}
    N(S') = N(S) + q.
    \end{equation}
    \item $A$ carries an \emph{energy increase bound} $d \in \mathbb{N}$ if for every basis state $|S\rangle$, every output state $|S'\rangle$ in the support of $A |S\rangle$ has excitation energy:
    \begin{equation}
    e(S') \le e(S) + d.
    \end{equation}
    \item The map is an \emph{exact charge-energy shift} $(q,d)\in\mathbb Z^2$ if every output configuration in $A|S\rangle$ satisfies both $N(S')=N(S)+q$ and $e(S')=e(S)+d$.
\end{enumerate}
\end{definition}

\begin{theorem}[Filtered Operator Action]
\label{thm:filtered_action}
Let $A$ be a linear operator with charge shift $q \in \mathbb{Z}$ and energy increase bound $d \in \mathbb{N}$.
Then $A$ maps the budget subspace $B(N, K)$ into $B(N+q, K+d)$:
\begin{equation}
A\big( B(N, K) \big) \subseteq B(N+q, K+d).
\end{equation}
\end{theorem}

\begin{proof}
Every vector in $B(N, K)$ is a linear combination of basis vectors $|S\rangle$ satisfying $N(S) = N$ and $e(S) \le K$. Applying $A$, each output basis vector $|S'\rangle$ satisfies $N(S') = N + q$ and $e(S') \le e(S) + d \le K + d$, which lies in $B(N+q, K+d)$.
\end{proof}

\begin{theorem}[Unified Margin Inequality]
\label{thm:unified_margin}
Let density operators with maximum mode cutoff $M \ge 1$ act on states in the charge-box budget $\mathcal{B}_{K, N_{\max}}$. If intermediate states undergo cumulative upward energy excursions of at most $K_{\mathrm{excursion}}$, the condition ensuring that all active modes stay strictly within the linear interior of the band without hitting the Nyquist boundaries is:
\begin{equation}
2M + K + K_{\mathrm{excursion}} + N_{\max} \le h.
\end{equation}
Under this condition, no density shift triggers an Umklapp boundary wrap, so the specified mode transfers retain their unwrapped integer action on the stated budget.
\end{theorem}

\begin{proof}
For every intermediate state, any highest occupied mode is at most $N+K+K_{\mathrm{excursion}}$, while any lowest unoccupied mode is at least $N-K-K_{\mathrm{excursion}}+1$. Using $|N|\le N_{\max}$ and allowing transfers up to $M$ with the corresponding edge buffer, the unified condition $2M+K+K_{\mathrm{excursion}}+N_{\max}\le h$ keeps both extremal modes in the linear interior and prevents Umklapp wrapping.
\end{proof}

\begin{remark}[Finite Algebraic Foundations]
The rank formula, coordinate projectors, filtered maps, and word excursion bounds follow from finite sums on the occupation basis. The vacuum projectors and displacement description make the energy-budget regularization explicit without invoking an abstract dimension count or an infinite-volume limit.
\end{remark}

\section{Word Excursion Tracking and Compression Remainders}

For a product of operators $W=A_1A_2\cdots A_s$ with exact energy shifts $\Delta e_j$ and charge shifts $q_j$, the rightmost factor acts first. Define its maximum upward energy excursion by
\begin{equation}
X(W)=\max\left(0,\max_{1\le r\le s}\sum_{j=1}^{r}\Delta e_{s-j+1}\right).
\end{equation}

\begin{theorem}[Word Evaluation Safety]
\label{thm:word_evaluation_safety}
Suppose $\psi\in B(N,K)$ and $K+X(W)\le K_{\max}$. After applying the first $r$ factors from the right, every intermediate vector lies in
\begin{equation}
B\left(N+\sum_{j=1}^{r}q_{s-j+1},K_{\max}\right),
\end{equation}
so every intermediate state remains within the energy budget $K_{\max}$.
\end{theorem}

\begin{proof}
After $r$ right-to-left applications the cumulative energy shift is $\sum_{j=1}^{r}\Delta e_{s-j+1}\le X(W)$. Starting from energy at most $K$, the intermediate energy is therefore at most $K+X(W)\le K_{\max}$; the charge shift is the sum of the corresponding $r$ charge shifts.
\end{proof}

\begin{proposition}[Compression Remainder Identity]
\label{prop:compression_remainder}
For a coordinate projector $P=P_U$ and operators $A,B\in\operatorname{End}_{\mathbb C}(\mathcal F(L))$, the product of compressed operators differs from the compressed product by the exact intermediate leakage term
\begin{equation}
PABP-(PAP)(PBP)=PA(I-P)BP.
\end{equation}
\end{proposition}

\begin{proof}
Expanding $PABP-PAPBP$ and factoring on the left and right gives $PA(I-P)BP$.
\end{proof}

\section{Chapter Summary and Literature Comparison}

In this chapter, we developed the Energy Budget regularization scheme:
\begin{itemize}
    \item Admissible charge sectors $N \in \{-h, \dots, h\}$ possess unique ground configurations $\mathcal{S}_N = \{-h+1, \dots, N\}$.
    \item The sector ground reference offset $E_0(N) = N(N+1)/2$ evaluates the macroscopic charging energy in exact integer arithmetic.
    \item Defined particle displacements $d_i = s_i - g_i$ and proved the Rank Formula $e(S) = \sum d_i \ge 0$, establishing ground state uniqueness $e(S) = 0 \iff S = \mathcal{S}_N$.
    \item Connected ground state non-degeneracy to the gapless $\mathcal{O}(1/L)$ critical scaling of the 1D Fermi gas.
    \item Proved the Frozen Margins Theorem: states with $e(S) \le K$ have rigidly filled deep seas ($k \le N - K$) and rigidly empty high modes ($k > N + K$).
    \item Defined budget subspaces $B(N, K) = \bigoplus_{E=0}^K \mathcal{H}(N, E)$, specified filtered operator parameters $(q, d)$, and proved the Unified Margin Inequality $2M + K + K_{\mathrm{excursion}} + N_{\max} \le h$.
    \item Established the coordinate-projector algebra, tracked cumulative energy excursions of operator words, and identified the exact leakage term produced by compressing an operator product.
\end{itemize}

\begin{quote}
\item[]\paragraph{Comparison with Standard Continuum Literature.}
In standard continuum bosonization (Giamarchi~\cite{giamarchi2004}, Gogolin et al.~\cite{gogolin1998}, and von Delft and Schoeller~\cite{vondelft1998}), many treatments use smooth regulators, such as exponential damping factors $e^{-\alpha |k|}$ or Gaussian filters, to control momentum sums. The present construction instead specifies finite-dimensional orthogonal subspaces directly through integer charge and excitation-energy conditions.

The energy budget $B(N,K)$ is an exact finite-dimensional orthogonal subspace. Theorem~\ref{thm:unified_margin} gives a sufficient margin under which the specified finite-lattice mode transfers avoid boundary wrapping.
\end{quote}
